\begin{abstract}
大语言模型（LLM）直接“读年报、写研报”有三类在投研场景中不可接受的失败：编造或记错财务数字；在增速、占比与估值上心算出错；用参数记忆补写没有出处的外部信息。此外，研报质量本身缺少可复核的客观评价，使“改进”难以度量。本文介绍 \sysname{}，一个面向 A 股上市公司的可追溯投研多智能体系统，覆盖财报分析、估值建模与研报纠错三项任务，并能把一份年报加工为 25～32 页、60 个以上图表的卖方体例深度研报和公式联动的估值模型。

系统的核心设计是\textbf{理解与计算分离}：模型负责读、抄、写与判断，确定性程序负责校验、计算与拦截。具体包括：（1）证据约束的双通道抽取，每个数字带页码级来源并逐字回查原文，规则通道与模型通道冲突时不静默取舍；（2）单一真值表驱动的分析与估值链（分部量价 → 三表联动 → FCFF 折现 → 股权价值桥 → 程序评级）；（3）按展示精度设定容差的\textbf{数字护栏}与约束改写的\textbf{数字锁}；（4）由五个上下文隔离的无头会话和 11 道程序闸门组成的研报流水线，先搜集外部资料、后分析写作；（5）四层评分标准与盲评闭环，评分标准对生成方不可见，并用重复评审、程序制造缺陷与人工复核三组证据校准评审模型本身。

在三家公司的真实年报上，314 个人工核验数值中规则通道覆盖 293 个，错配为 0；流水线迭代后，两家公司深度研报的质检得分由 74、67 分提升到 90、88 分（同一评分标准、单次盲评，非对照实验）；评审模型重复评分的逐条一致率为 86\%～91\%，判出问题经人工确认的精确率为 84\%。所有数值计算均可在保存的输入上离线复算。
\end{abstract}
